La auditoría de sesgo observacional se incorporó porque los exoplanetas no se descubren de manera uniforme. Cada método de detección favorece regiones distintas del espacio planetario. Por ejemplo, el método de tránsito detecta con mayor facilidad planetas de periodo corto, cercanos a su estrella, grandes y alineados con nuestra línea de visión. La velocidad radial favorece planetas masivos o de periodo moderado alrededor de estrellas donde la señal gravitacional sea medible. La imagen directa favorece planetas gigantes, jóvenes, calientes, luminosos y separados de su estrella. Por esta razón, una región Mapper puede aparecer no solamente porque los planetas sean físicamente parecidos, sino porque fueron detectados mediante un mismo procedimiento observacional.

En este reporte, las variables \texttt{discoverymethod}, \texttt{disc\_year} y \texttt{disc\_facility} se usan como metadata externa para auditar sesgo observacional. No son variables de entrada de Mapper. Esta decisión es metodológicamente importante: \texttt{discoverymethod} no describe qué es físicamente el planeta, sino cómo fue encontrado. Si se usara como coordenada del Mapper, podría contaminar la geometría del grafo y agrupar planetas por historia observacional en lugar de por propiedades físicas, orbitales o térmicas.

La pregunta de esta auditoría es entonces:

\begin{quote}
¿Las regiones del grafo Mapper están alineadas con propiedades físicas u orbitales, o están enriquecidas por el proceso de descubrimiento?
\end{quote}

Para responderla, se calculó la composición de cada nodo por método de descubrimiento. Para un nodo \(v\) con conjunto de miembros \(S_v\), la distribución por método se define como:

\[
p_v(m)=\frac{|\{i\in S_v : d_i=m\}|}{|S_v|},
\]

donde \(d_i\) es el método de descubrimiento del planeta \(i\), y \(m\) recorre los métodos observados en el catálogo. A partir de esta distribución se calculó la pureza por método:

\[
P_v = \max_m p_v(m),
\]

que mide la fracción del método dominante dentro del nodo. Si \(P_v=1\), todos los planetas del nodo fueron descubiertos por el mismo método. Si \(P_v\) es menor, el nodo mezcla varios métodos.

También se calculó la entropía por método:

\[
H_v=-\sum_m p_v(m)\log p_v(m).
\]

La entropía mide qué tan mezclados están los métodos de descubrimiento dentro del nodo. Una entropía baja indica dominancia de un método; una entropía alta indica una composición más diversa. Por lo tanto, un nodo con pureza alta y entropía baja es observacionalmente más sospechoso que un nodo con mezcla de métodos.

Además de pureza y entropía, se evaluó la divergencia entre la composición local del nodo y la composición global del catálogo. Esta comparación es necesaria porque algunos métodos, como tránsito, pueden ser frecuentes en el catálogo completo. Un nodo con muchos planetas de tránsito no es necesariamente sorprendente si el catálogo global ya está dominado por tránsito. En cambio, un nodo cuya composición por método se aleja fuertemente de la distribución global muestra enriquecimiento local por método de descubrimiento.

Las variables \texttt{disc\_year} y \texttt{disc\_facility} añaden una capa temporal e instrumental a la auditoría. \texttt{disc\_year} permite detectar si una región está concentrada en una época de descubrimiento, mientras que \texttt{disc\_facility} permite identificar regiones dominadas por una misión, telescopio o instalación específica. Si un nodo está dominado simultáneamente por un método, una facilidad y una ventana temporal estrecha, la evidencia de sesgo observacional se vuelve más fuerte.

Una región se considera observacionalmente sospechosa cuando presenta una o varias de las siguientes señales:

\begin{itemize}
    \item alta pureza por \texttt{discoverymethod};
    \item baja entropía por método;
    \item alta divergencia respecto a la distribución global del catálogo;
    \item dominancia de una sola \texttt{disc\_facility};
    \item concentración en una ventana estrecha de \texttt{disc\_year};
    \item enriquecimiento significativo en la prueba de permutación.
\end{itemize}

Es importante subrayar que esta asociación no prueba causalidad. Un nodo dominado por \texttt{Transit}, \texttt{Radial Velocity} o \texttt{Imaging} no demuestra por sí mismo que el método haya causado la estructura. Sin embargo, sí indica que la región debe interpretarse con cautela, porque puede estar mezclando señal astrofísica con selección observacional.

\subsection*{Nodos enriquecidos por método de descubrimiento}

La auditoría identificó nodos con enriquecimiento fuerte por \texttt{discoverymethod}. Los casos más extremos aparecen en los grafos conjuntos y están dominados por \texttt{Imaging}. Por ejemplo, varios nodos en \texttt{joint\_pca2} y \texttt{joint\_no\_density\_pca2} presentan fracción dominante igual a 1.000, valores \(Z\) muy altos y \(p=0.001\). Estos nodos son especialmente sospechosos porque \texttt{Imaging} no observa uniformemente toda la población planetaria: tiende a favorecer planetas gigantes, jóvenes, luminosos y separados de su estrella.

\begin{table}[H]
\centering
\scriptsize
\caption{Nodos con mayor enriquecimiento por metodo de descubrimiento.}
\label{tab:top_enriched_nodes}
\begin{adjustbox}{max width=\linewidth}
\begin{tabular}{lllrrrr}
\toprule
config\_id & node\_id & dominant\_discoverymethod & observed\_dominant\_method\_fraction & enrichment\_z & empirical\_p\_value & n\_members \\
\midrule
joint\_pca2\_cubes10\_overlap0p35 & cube61\_cluster0 & Imaging & 1.000 & 46.721 & 0.001 & 8 \\
joint\_no\_density\_pca2\_cubes10\_overlap0p35 & cube61\_cluster0 & Imaging & 1.000 & 46.556 & 0.001 & 10 \\
joint\_pca2\_cubes10\_overlap0p35 & cube62\_cluster0 & Imaging & 1.000 & 39.721 & 0.001 & 5 \\
joint\_pca2\_cubes10\_overlap0p35 & cube47\_cluster0 & Imaging & 1.000 & 39.001 & 0.001 & 7 \\
joint\_pca2\_cubes10\_overlap0p35 & cube54\_cluster0 & Imaging & 1.000 & 39.001 & 0.001 & 7 \\
joint\_pca2\_cubes10\_overlap0p35 & cube48\_cluster1 & Imaging & 1.000 & 34.440 & 0.001 & 6 \\
joint\_pca2\_cubes10\_overlap0p35 & cube55\_cluster0 & Imaging & 1.000 & 34.440 & 0.001 & 6 \\
orbital\_pca2\_cubes10\_overlap0p35 & cube6\_cluster0 & Transit & 0.928 & 32.343 & 0.001 & 3316 \\
joint\_no\_density\_pca2\_cubes10\_overlap0p35 & cube61\_cluster1 & Imaging & 1.000 & 29.951 & 0.001 & 4 \\
joint\_no\_density\_pca2\_cubes10\_overlap0p35 & cube62\_cluster0 & Imaging & 1.000 & 29.951 & 0.001 & 4 \\
joint\_no\_density\_pca2\_cubes10\_overlap0p35 & cube68\_cluster0 & Imaging & 1.000 & 29.951 & 0.001 & 4 \\
thermal\_pca2\_cubes10\_overlap0p35 & cube26\_cluster0 & Radial Velocity & 0.586 & 29.905 & 0.001 & 720 \\
joint\_no\_density\_pca2\_cubes10\_overlap0p35 & cube50\_cluster0 & Imaging & 1.000 & 29.623 & 0.001 & 5 \\
joint\_no\_density\_pca2\_cubes10\_overlap0p35 & cube51\_cluster0 & Imaging & 1.000 & 29.623 & 0.001 & 5 \\
joint\_no\_density\_pca2\_cubes10\_overlap0p35 & cube59\_cluster0 & Imaging & 1.000 & 29.623 & 0.001 & 5 \\
\bottomrule
\end{tabular}
\end{adjustbox}
\vspace{0.4em}\begin{minipage}{0.96\linewidth}\footnotesize Se muestran 15 de 496 filas; el CSV fuente contiene la tabla completa.\end{minipage}
\end{table}


La Tabla anterior debe leerse como una lista de regiones donde la topología local se alinea fuertemente con el método de descubrimiento. En particular, los nodos dominados por \texttt{Imaging} en los espacios conjuntos no deben interpretarse directamente como una población física limpia. Pueden ser científicamente interesantes, porque corresponden a planetas con propiedades extremas o poco comunes, pero su composición está fuertemente condicionada por el método que permitió detectarlos.

En el grafo orbital también aparecen nodos relevantes. Algunos nodos están dominados por \texttt{Transit}, como \texttt{cube6\_cluster0}, con fracción dominante de 0.928, y \texttt{cube12\_cluster0}, con fracción dominante de 0.877. Estos nodos son importantes porque el grafo orbital usa \texttt{pl\_orbper} y \texttt{pl\_orbsmax}; justamente las variables orbitales pueden estar muy ligadas al sesgo de tránsito, que favorece planetas cercanos y de periodo corto. Por lo tanto, una región orbital dominada por \texttt{Transit} puede reflejar una combinación de arquitectura orbital real y sesgo de detección.

También aparecen casos orbitales dominados por \texttt{Radial Velocity}, como \texttt{cube25\_cluster0}, con fracción dominante de 0.605, y nodos pequeños con presencia de \texttt{Imaging}, como \texttt{cube32\_cluster5}. Estos ejemplos no tienen la misma lectura que los nodos grandes dominados por tránsito: su tamaño, imputación y composición deben revisarse antes de asignarles una interpretación fuerte. En particular, los nodos pequeños pueden mostrar pureza alta simplemente por bajo número de miembros, por lo que deben leerse como señales de auditoría y no como evidencia concluyente.

\subsection*{Componentes enriquecidos}

La auditoría también se realizó a nivel de componentes conectadas. Esta escala es útil porque un componente puede reunir varios nodos relacionados y mostrar si una porción completa del grafo está dominada por un método de descubrimiento. Los componentes más sospechosos son aquellos pequeños o medianos dominados completamente por un solo método.

\begin{table}[H]
\centering
\scriptsize
\caption{Componentes con mayor concentracion observacional.}
\label{tab:component_discovery_bias}
\begin{adjustbox}{max width=\linewidth}
\begin{tabular}{lrrlrrlr}
\toprule
config\_id & component\_id & n\_members & dominant\_discoverymethod & dominant\_discoverymethod\_fraction & discoverymethod\_entropy & dominant\_disc\_facility & disc\_year\_median \\
\midrule
joint\_no\_density\_pca2\_cubes10\_overlap0p35 & 6 & 5 & Imaging & 1.000 & -0.000 & European Space Agency (ESA) Gaia Satellite & 2019.000 \\
joint\_no\_density\_pca2\_cubes10\_overlap0p35 & 7 & 10 & Imaging & 1.000 & -0.000 & W. M. Keck Observatory & 2018.000 \\
joint\_no\_density\_pca2\_cubes10\_overlap0p35 & 8 & 4 & Imaging & 1.000 & -0.000 & Paranal Observatory & 2016.500 \\
joint\_pca2\_cubes10\_overlap0p35 & 5 & 7 & Imaging & 1.000 & -0.000 & European Space Agency (ESA) Gaia Satellite & 2019.000 \\
joint\_pca2\_cubes10\_overlap0p35 & 6 & 8 & Imaging & 1.000 & -0.000 & Paranal Observatory & 2018.000 \\
thermal\_pca2\_cubes10\_overlap0p35 & 12 & 5 & Imaging & 0.600 & 1.371 & Multiple Observatories & 2023.000 \\
orbital\_pca2\_cubes10\_overlap0p35 & 18 & 4 & Microlensing & 0.500 & 1.500 & KMTNet & 2020.500 \\
thermal\_pca2\_cubes10\_overlap0p35 & 13 & 4 & Radial Velocity & 0.750 & 0.811 & European Space Agency (ESA) Gaia Satellite & 2015.000 \\
orbital\_pca2\_cubes10\_overlap0p35 & 23 & 4 & Imaging & 0.500 & 1.500 & Hubble Space Telescope & 2014.500 \\
thermal\_pca2\_cubes10\_overlap0p35 & 6 & 8 & Radial Velocity & 0.875 & 0.544 & La Silla Observatory & 2014.000 \\
thermal\_pca2\_cubes10\_overlap0p35 & 15 & 6 & Microlensing & 0.667 & 0.918 & KMTNet & 2023.000 \\
orbital\_pca2\_cubes10\_overlap0p35 & 13 & 5 & Radial Velocity & 0.800 & 0.722 & Multiple Observatories & 2014.000 \\
thermal\_pca2\_cubes10\_overlap0p35 & 3 & 10 & Radial Velocity & 0.900 & 0.469 & Multiple Observatories & 2010.000 \\
orbital\_pca2\_cubes10\_overlap0p35 & 12 & 5 & Radial Velocity & 1.000 & -0.000 & La Silla Observatory & 2011.000 \\
orbital\_pca2\_cubes10\_overlap0p35 & 14 & 4 & Radial Velocity & 1.000 & -0.000 & W. M. Keck Observatory & 2017.000 \\
\bottomrule
\end{tabular}
\end{adjustbox}
\vspace{0.4em}\begin{minipage}{0.96\linewidth}\footnotesize Se muestran 15 de 110 filas; el CSV fuente contiene la tabla completa.\end{minipage}
\end{table}


En los resultados, varios componentes de \texttt{joint\_no\_density\_pca2} y \texttt{joint\_pca2} aparecen completamente dominados por \texttt{Imaging}. Algunos de ellos tienen además imputación moderada o alta, por ejemplo componentes con imputación de 0.333, 0.500, 0.286 o 0.429. Esta combinación es importante: cuando una región está dominada por un método de descubrimiento y además depende de imputación, su interpretación física se debilita. En esos casos, la región debe clasificarse con mayor probabilidad como observacional o mixta.

En el grafo orbital también pueden existir componentes dominados por \texttt{Transit} o \texttt{Radial Velocity}. Sin embargo, varios de esos componentes son pequeños, por lo que funcionan mejor como alertas de auditoría que como conclusiones fuertes. Esta diferencia muestra por qué la auditoría debe combinar varias métricas: pureza, entropía, divergencia, tamaño del nodo o componente, imputación y contexto astrofísico.

\subsection*{Lectura metodológica de la auditoría}

La conclusión principal de esta auditoría es que varias regiones Mapper están significativamente enriquecidas por método de descubrimiento. Esto no invalida los grafos, pero sí cambia su interpretación. Una estructura topológica puede ser informativa y, al mismo tiempo, estar parcialmente asociada con la manera en que los planetas fueron detectados.

En particular, las regiones conjuntas dominadas por \texttt{Imaging} y las regiones orbitales dominadas por \texttt{Transit} deben interpretarse como evidencia mixta: potencialmente astrofísica, pero también observacionalmente sesgada. Esta lectura es más fuerte que simplemente aceptar o rechazar el grafo. Permite separar regiones con señal física más limpia de regiones donde la topología observada puede estar influida por el catálogo.

Por lo tanto, la auditoría de sesgo observacional funciona como un filtro de interpretación. Una región con coherencia física, baja imputación y baja concentración por método puede considerarse una candidata más robusta. Una región con coherencia física pero alta dominancia por \texttt{discoverymethod} debe leerse como mixta. Una región dominada por método, facilidad o época, y además con alta imputación, debe leerse como observacional o débil.

Esta etapa es central para el objetivo del proyecto: no basta con encontrar estructura en Mapper; hay que evaluar si esa estructura se sostiene como señal astrofísica o si refleja el proceso de observación. Bajo esta lectura, el reporte no propone una taxonomía final de exoplanetas, sino una clasificación crítica de regiones candidatas según la evidencia que las sostiene.